\chapter{镜像对称的构造}\label{chap:mirror-constructions}
\SourcePage{76}

本章介绍镜像流形的各种构造。其中大部分构造涉及环面簇，但并非全部如此。特别地，我们将从环面几何的角度解释五次超曲面的镜像对称。

\section{Batyrev 镜像构造}\label{sec:batyrev-mirror}

\hyperref[chap:introduction]{第 1 章}从超共形场论的角度讨论了镜像对称为何合理，却没有说明如何寻找给定 Calabi--Yau 三维簇的镜像。物理中最早的镜像流形显式例子由 Greene 与 Plesser \SourceCite{GPl} 给出，他们将镜像构造成加权射影空间中 Fermat 超曲面的商。\hyperref[chap:quintic]{第 2 章}给出的五次三维簇的镜像，就是这种构造的一个例子。Roan 观察到，环面几何提供了描述这类镜像的自然框架 \SourceCite{Roan1}。随后，Batyrev 找到了一种漂亮的环面几何描述，大大推广了 Greene--Plesser 构造。

Batyrev 的构造关键地使用了\hyperref[sec:fano-reflexive]{第 3.5 节}介绍的自反多胞形。那里已看到，自反多胞形 $\Delta$ 对应于 Fano 环面簇，而且 $\Delta$ 有一个同样自反的对偶 $\Delta^\circ$。本节将解释，$\Delta$ 与 $\Delta^\circ$ 之间的对偶性，如何导出某些与 $\mathbb P_\Delta$、$\mathbb P_{\Delta^\circ}$ 密切相关的环面簇中 Calabi--Yau 超曲面族之间的对偶性。\index{Batyrev镜像构造@Batyrev 镜像构造|hyperpage}

\subsection{环面簇中的 Calabi--Yau 超曲面}\label{subsec:toric-cy-hypersurfaces}

Batyrev 构造的第一步，是观察到环面簇中的 Calabi--Yau 超曲面自然来自自反多胞形。自反多胞形 $\Delta$ 给出环面簇 $\mathbb P_\Delta$。不过，对于镜像对称，$\mathbb P_\Delta$ 的奇点可能过于严重。我们需要能在其上进行 Hodge 理论的对象，因此将使用细化 $\mathbb P_\Delta$ 的扇的单纯扇。回顾，$\mathbb P_\Delta$ 的扇是 $\Delta$ 的法扇，由向 $\Delta^\circ\subset N_{\mathbb R}$ 的各真面作出的锥组成。此外，因为 $\Delta$ 自反，法扇的一维锥生成元就是 $\Delta^\circ$ 的各顶点，它们属于 $\Delta^\circ\cap N-\{0\}$。我们将考虑 $N_{\mathbb R}$ 中的以下扇。

\begin{definition}\label{def:projective-subdivision}
给定自反多胞形 $\Delta\subset M_{\mathbb R}$，若 $N_{\mathbb R}$ 中的扇 $\Sigma$ 具有以下性质，则称它为\emph{射影细分}：
\begin{itemize}
\item $\Sigma$ 细化 $\Delta$ 的法扇；
\item $\Sigma(1)\subset\Delta^\circ\cap N-\{0\}$；
\item $X_\Sigma$ 射影且单纯。
\end{itemize}
此外，若 $\Sigma(1)=\Delta^\circ\cap N-\{0\}$，则称 $\Sigma$ 为\emph{极大}的。
\end{definition}

可以证明，极大射影细分 $\Sigma$ 存在 \SourceCite{OdP}。这样的 $\Sigma$ 有多种选择。如\hyperref[sec:gkz-decomposition]{第 3.4 节}所述，每种选择给出
\SourcePage{77}
GKZ 分解中的一个不同的锥，也给出相关规范线性 $\sigma$ 模型中的一个不同的相，见\hyperref[app:physical-theories]{附录 B}。\hyperref[subsec:multiple-mirrors]{第 4.1.4 小节}将更详细地讨论这一点。

给定射影细分 $\Sigma$，得到环面簇 $X_\Sigma$ 以及双有理态射 $f:X_\Sigma\to\mathbb P_\Delta$，它们具有以下良好性质。
\begin{lemma}\label{lem:projective-subdivision-properties}
若 $\Sigma$ 是射影细分，则：
\begin{enumerate}[label=(\roman*)]
\item $X_\Sigma$ 是 Gorenstein 轨形；
\item $\Delta$ 是 $-K_{X_\Sigma}$ 对应的多胞形；
\item $-K_{X_\Sigma}$ 是半丰沛的，并且是大的；具体说，$-K_{X_\Sigma}$ 由整体截面生成，且 $(-K_{X_\Sigma})^n>0$；
\item 映射 $f:X_\Sigma\to\mathbb P_\Delta$ 是保典范的，即 $f^*(K_{\mathbb P_\Delta})=K_{X_\Sigma}$。
\end{enumerate}
此外，若 $\Sigma$ 是极大射影细分，则 $X_\Sigma$ 具有终端奇点。
\end{lemma}
\begin{proof}
因为 $\Delta^\circ$ 自反，它的每个刻面 $\Gamma^\circ$ 都由 $\langle m,v\rangle=-1$ 定义，其中某个 $m=m_{\Gamma^\circ}\in M$。于是当 $v_\rho\in\Gamma^\circ$ 时，有 $\langle m,v_\rho\rangle=-1$，根据\hyperref[sec:cones-fans]{第 3.1 节}的准则，容易推出 $-K_{X_\Sigma}=\sum_\rho D_\rho$ 是 Cartier 除子。因此 $X=X_\Sigma$ 是 Gorenstein 的；又因 $\Sigma$ 单纯，它也是轨形。此外，等式 $\langle m,v_\rho\rangle=-1$ 表明，与除子 $-K_X$ 对应的多胞形恰好是 $\Delta$。由对偶性，各 $m$ 就是 $\Delta$ 的顶点，\hyperref[subsec:toric-polytopes]{第 3.2.1 小节}表明 $-K_X$ 由整体截面生成。最后，由 \SourceCite[第 5.3 节]{Fulton3}，交数 $(-K_{X_\Sigma})^n=n!\operatorname{vol}(\Delta)$；由于 $\Delta$ 是 $n$ 维的，这个数为正。

为证明引理余下的部分，研究向 $\Sigma$ 加入新的一维锥生成元如何影响 $X$ 的典范类。回顾定义~\ref{def:fano} 后的讨论，对任意环面簇，$K_X$ 是 Weil 除子 $-\sum_{\rho\in\Sigma(1)}D_\rho$。取锥 $\sigma\in\Sigma$，在其相对内部加入新的本原整生成元 $v'\in N$，所得环面簇记为 $X'$。已知 $\sigma$ 的生成元 $v_\rho$ 位于 $\Delta^\circ$ 的某个刻面 $\Gamma^\circ$ 上，因此 $\langle m,v_\rho\rangle=-1$，其中 $m=m_{\Gamma^\circ}$ 如上。关键计算是：若 $g:X'\to X$ 为自然映射，则
\begin{equation}\label{eq:toric-canonical-discrepancy}
K_{X'}=g^*(K_X)-\bigl(\langle m,v'\rangle+1\bigr)D',
\end{equation}
其中 $D'$ 是 $X'$ 上对应于 $v'$ 的除子。这个等式应理解为 Weil 除子的线性等价关系。例如，可用 \SourceCite[第 4 节]{Reid4} 的方法证明式~\eqref{eq:toric-canonical-discrepancy}。只需考察其中 $D'$ 的系数：在仿射开集 $X_\sigma\subset X$ 上，$K_X$ 与 $\operatorname{div}(\chi^m)$ 一致，因此 $g^*K_X$ 中 $D'$ 的系数为 $\langle m,v'\rangle$。

若 $v'\in\Delta^\circ\cap N-\{0\}$，则 $v'$ 位于刻面 $\Gamma^\circ$ 上，所以 $\langle m,v'\rangle=-1$，由式~\eqref{eq:toric-canonical-discrepancy} 立即得到 $K_{X'}=g^*(K_X)$。因为 $\Sigma(1)\subset\Delta^\circ\cap N$，$X_\Sigma$ 由 $\mathbb P_\Delta$ 通过依次加入 $\Delta^\circ\cap N$ 中的新顶点得到，容易推出 $K_{X_\Sigma}=f^*(K_{\mathbb P_\Delta})$。

另一方面，若 $\Sigma$ 极大，则 $\Sigma(1)=\Delta^\circ\cap N-\{0\}$。因此，任何新加入的本原整生成元 $v'$ 必定位于 $\Delta^\circ$ 之外；按自反的定义，$0$ 是唯一的内部格点。于是式~\eqref{eq:toric-canonical-discrepancy} 中 $\langle m,v'\rangle+1<0$。由\hyperref[app:singular-varieties]{附录 A} 的讨论，可见 $X$ 具有终端奇点。
\end{proof}

下面说明如何利用 $X_\Sigma$ 和 $\mathbb P_\Delta$ 上的反典范线性系构造 Calabi--Yau 簇。\index{射影细分|hyperpage}\index{保典范态射|hyperpage}

\SourcePage{78}
\begin{proposition}\label{prop:toric-cy-construction}
若 $\Delta$ 是 $n$ 维自反多胞形，则一般成员 $\overline V\in|-K_{\mathbb P_\Delta}|$ 是 $n-1$ 维 Calabi--Yau 簇。此外，若 $\Sigma$ 是射影细分，$X=X_\Sigma$，则：
\begin{enumerate}[label=(\roman*)]
\item 一般成员 $V\in|-K_X|$ 是 Calabi--Yau 轨形；
\item 若 $\Sigma$ 极大，则一般成员 $V\in|-K_X|$ 是极小 Calabi--Yau 轨形。
\end{enumerate}
\end{proposition}
\begin{proof}
先考虑 $\overline V$。根据\hyperref[sec:math-preview]{第 1.4 节}给出的 Calabi--Yau 簇的定义，必须证明 $\overline V$ 具有典范奇点、平凡的典范层，而且对 $0<k<n-1$，上同调 $H^k(\overline V,\mathcal O_{\overline V})=0$。

由于 Fano 环面簇 $\mathbb P_\Delta$ 是 Gorenstein 的，它至多具有典范奇点，见 \SourceCite{Batyrev4}。于是 Bertini 定理的一个版本保证，一般成员 $\overline V\in|-K_{\mathbb P_\Delta}|$ 也至多具有典范奇点，见 \SourceCite[定理 1.13]{Reid1}。还应注意
\[
\widehat\Omega_{\overline V}^{n-1}
\simeq\widehat\Omega_{\mathbb P_\Delta}^{n}
(-K_{\mathbb P_\Delta})\otimes\mathcal O_{\overline V}
\simeq\mathcal O_{\overline V},
\]
其中第一个同构是伴随公式；因为 $\mathbb P_\Delta$ 是 Cohen--Macaulay 的，且 $-K_{\mathbb P_\Delta}$ 是 Cartier 除子，该公式成立。最后一步是证明 $0<k<n-1$ 时 $H^k(\overline V,\mathcal O_{\overline V})=0$。由于
$\mathcal O_{\mathbb P_\Delta}(-\overline V)\simeq\mathcal O_{\mathbb P_\Delta}(K_{\mathbb P_\Delta})=\widehat\Omega_{\mathbb P_\Delta}^n$，得到正合列
\[
0\longrightarrow\widehat\Omega_{\mathbb P_\Delta}^n
\longrightarrow\mathcal O_{\mathbb P_\Delta}
\longrightarrow\mathcal O_{\overline V}\longrightarrow0,
\]
进而得到长正合列
\[
\cdots\longrightarrow H^k(\mathbb P_\Delta,\mathcal O_{\mathbb P_\Delta})
\longrightarrow H^k(\overline V,\mathcal O_{\overline V})
\longrightarrow H^{k+1}(\mathbb P_\Delta,\widehat\Omega_{\mathbb P_\Delta}^n)
\longrightarrow\cdots.
\]
然而，对 $k>0$ 有 $H^k(\mathbb P_\Delta,\mathcal O_{\mathbb P_\Delta})=0$；由 Serre--Grothendieck 对偶性，对 $k<n-1$ 还有
\[
H^{k+1}(\mathbb P_\Delta,\widehat\Omega_{\mathbb P_\Delta}^n)
\simeq H^{n-k-1}(\mathbb P_\Delta,\mathcal O_{\mathbb P_\Delta})^*=0.
\]
这给出所需的 $H^k(\overline V,\mathcal O_{\overline V})$ 的消失，因此 $\overline V$ 是 Calabi--Yau 簇。

现在令 $\Sigma$ 是射影细分，置 $X=X_\Sigma$。由引理~\ref{lem:projective-subdivision-properties}，$-K_X$ 半丰沛，所以线性系 $|-K_X|$ 没有基点。此外，$\Sigma$ 单纯，所以 $X$ 是轨形；将 Bertini 定理应用于定义 $X$ 的局部商的各不动点集，便保证一般成员 $V\in|-K_X|$ 是 $X$ 的子轨形。上面的全部论证仍然成立，因此 $V$ 是 Calabi--Yau 轨形。

最后设 $\Sigma$ 极大。根据定义~\ref{def:calabi-yau-variety}，$V$ 为极小 Calabi--Yau 簇的条件是它具有 Gorenstein、$\mathbb Q$ 析因、终端奇点。因为 $V$ 已是 Gorenstein 轨形，它的奇点自动为 Gorenstein 且 $\mathbb Q$ 析因。因此只需证明 $V$ 具有终端奇点。不过，由引理~\ref{lem:projective-subdivision-properties}，环境空间 $X$ 具有终端奇点，像 \SourceCite[定理 1.13]{Reid1} 的证明那样应用 Bertini 定理即可。
\end{proof}

由于 $f:X_\Sigma\to\mathbb P_\Delta$ 保典范，上面构造的 Calabi--Yau 超曲面 $V\subset X_\Sigma$ 是 $\mathbb P_\Delta$ 中一般反典范超曲面 $\overline V$ 经 $f$ 的严格变换。还可以证明，诱导映射 $V\to\overline V$ 也保典范。

当 $\Sigma$ 是极大射影细分时，簇 $V\subset X_\Sigma$ 是 $\overline V\subset\mathbb P_\Delta$ 在 Mori 理论意义下的极小模型。Batyrev 称 $V$ 为 $\overline V$ 的\emph{极大射影保典范部分消解}，简称 \emph{MPCP 消解}。

当 $V$ 是三维簇时，情形尤其好。由\hyperref[app:singular-varieties]{附录 A}，三维 Gorenstein 轨形的终端奇点实际上是光滑点。因此，对极大射影细分，命题~\ref{prop:toric-cy-construction} 构造的极小 Calabi--Yau 三维簇 $V\subset X_\Sigma$ 都是光滑的。\index{MPCP消解@MPCP 消解|hyperpage}

\SourcePage{79}
\subsection{环面 Calabi--Yau 超曲面的镜像}\label{subsec:batyrev-pair}

现在可以描述 Batyrev 镜像构造。固定自反多胞形 $\Delta$，令 $\Sigma$ 是定义~\ref{def:projective-subdivision} 中其法扇的一个极大射影细分。这样得到由 $V\in|-K_{X_\Sigma}|$ 给出的 Calabi--Yau 超曲面族 $V\subset X_\Sigma$。如\hyperref[chap:introduction]{第 1 章}所述，这个族中的每个成员都带有一个复化 Kähler 锥，从而得到族 $\{(V,\omega)\}$，它由 Calabi--Yau 超曲面及其复化 Kähler 类的一切可能配对组成。

不过，$\Delta$ 自反意味着 $\Delta^\circ$ 自反，因此可对 $\Delta^\circ$ 重复上述构造。对 $\Delta^\circ$ 在 $M_{\mathbb R}$ 中的法扇的一个极大射影细分 $\Sigma^\circ$，$|-K_{X_{\Sigma^\circ}}|$ 给出 Calabi--Yau 超曲面族 $V^\circ\subset X_{\Sigma^\circ}$。加入各复化 Kähler 锥之后，得到族 $\{(V^\circ,\omega^\circ)\}$，称为原族 $\{(V,\omega)\}$ 的 \emph{Batyrev 镜像}。

物理学认为，这两个族是物理学定义~\ref{def:physics-mirror-pair} 意义下的镜像对。更精确地说，有以下猜想。
\begin{physicsconjecture}\label{conj:batyrev-physics}
族 $\{(V,\omega)\}$ 与 $\{(V^\circ,\omega^\circ)\}$ 诱导同构的超共形场论；它们的 $N=2$ 超共形表示除\hyperref[sec:physics]{第 1.1 节}讨论的符号改变外完全相同。
\end{physicsconjecture}

另一种说法是，族 $\{(V,\omega)\}$ 与 $\{(V^\circ,\omega^\circ)\}$ 分别表示 $N=2$ 超共形场论（SCFT）模空间\emph{同一个}连通分支中的开子集。严格说来，虽然 Batyrev 镜像构造适用于任意维数的自反多胞形，这个猜想只适用于 $V$ 与 $V^\circ$ 为三维簇的情形，此时 $\Delta$ 是四维的。

物理中尚未证明这个镜像对称猜想，不过 \SourceCite{MP2} 作过尝试。在某些加权射影空间中具有适当次数的 Fermat 超曲面这一特殊情形，Greene--Plesser 构造 \SourceCite{GPl} 从物理上证明了 $\{(V,\omega)\}$ 与 $\{(V^\circ,\omega^\circ)\}$ 构成镜像对。一般情形仍未解决。

由于 SCFT 缺少严格基础，物理学猜想~\ref{conj:batyrev-physics} 并非通常数学意义下的猜想。不过，可以把这种 SCFT 相等关系的数学推论精确地表述出来。特别地，\hyperref[chap:introduction]{第 1 章}已看到，镜像对称意味着下列断言应当成立：
\begin{itemize}
\item 上同调群 $H^{1,1}(V)$，相应地 $H^{1,1}(V^\circ)$，应同构于 $H^{n-2,1}(V^\circ)$，相应地 $H^{n-2,1}(V)$。回顾，$V$ 与 $V^\circ$ 的维数为 $n-1$，其中 $n=\dim\Delta$。
\item 根据镜像对称的图~\eqref{eq:mirror-foliations}，$V$，相应地 $V^\circ$，的 Kähler 模应通过镜像映射局部同构于 $V^\circ$，相应地 $V$，的复模。
\item $V$，相应地 $V^\circ$，的 A 模型关联函数应通过镜像映射，对应于经适当归一化的 $V^\circ$，相应地 $V$，的 B 模型关联函数。
\end{itemize}
数学家面临的挑战是证明这些推论。当 $V^\circ$ 是 $V$ 的 Batyrev 镜像时，已有实质性进展。\hyperref[subsec:batyrev-hodge]{第 4.1.3 小节}将解决第一项；\hyperref[chap:moduli]{第~\ref*{chap:moduli} 章}研究环面 Calabi--Yau 超曲面的复模与 Kähler 模时，将处理第二项。第三项则要等到\hyperref[chap:quantum-cohomology]{第~\ref*{chap:quantum-cohomology} 章}，那里将定义\IndexLocation{indexloc:0748:79}数学镜像对，并提出 Hodge 理论的\IndexLocation{indexloc:0795:79}环面镜像猜想。\hyperref[chap:mirror-theorem]{第~\ref*{chap:mirror-theorem} 章}将讨论若干特殊情形的证明。

\SourcePage{80}
\subsection{Batyrev 镜像的 Hodge 数}\label{subsec:batyrev-hodge}

要证明两个有限维向量空间同构，只须证明它们的维数相同。因此，上面的第一项要求比较环面 Calabi--Yau 超曲面 $V$ 与其 Batyrev 镜像 $V^\circ$ 的 Hodge 数。

仍假定 $\Delta$ 是 $n$ 维自反多胞形，$V$ 是 $X_\Sigma$ 中的一般反典范超曲面，其中 $\Sigma$ 是 $\Delta$ 的法扇的极大射影细分。对 $\Delta^\circ$ 应用这个过程，得到 Batyrev 镜像族 $V^\circ$。\SourceCite{Batyrev4,AGM1} 证明，$V$ 与 $V^\circ$ 的 Hodge 数有如下关系。
\begin{theorem}\label{thm:batyrev-hodge}
若 $V^\circ$ 是 $V$ 的 Batyrev 镜像，则
\[
h^{1,1}(V)=h^{n-2,1}(V^\circ),\qquad
h^{n-2,1}(V)=h^{1,1}(V^\circ).
\]
\end{theorem}
\begin{proof}
当 $n=2$ 或 $3$ 时，结论分别由椭圆曲线与 K3 曲面的 Hodge 数得到。以下设 $n\geq4$。
先定义 $H^{1,1}(V)$ 与 $H^{n-2,1}(V)$ 中的一些有趣的子空间。对每个 $v_i\in\Sigma(1)$，令 $D_i\subset X_\Sigma$ 为相应的 $T_N$ 不变除子。这些除子在超曲面 $V\subset X_\Sigma$ 上限制为除子 $D_i'$；其上同调类生成一个复数域上的子空间，记为
\begin{equation}\label{eq:toric-hodge-subspace}
H_{\mathrm{toric}}^{1,1}(V)\subset H^{1,1}(V).
\end{equation}
再考虑 $H^{n-2,1}(V)\simeq H^1(V,T_V)$。首先注意，虽然 $V$ 可能有奇点，它却是至多具有终端奇点的 Gorenstein 轨形；因此，根据命题~\ref{prop:terminal-gorenstein-orbifold-deformations}，$H^1(V,T_V)$ 分类 $V$ 的无穷小形变。由超曲面 $V\in|-K_{X_\Sigma}|$ 确定的无穷小形变给出子空间，记为
\begin{equation}\label{eq:polynomial-hodge-subspace}
H_{\mathrm{poly}}^{n-2,1}(V)\subset H^{n-2,1}(V).
\end{equation}
我们断言，这些子空间的维数可用 $\Delta$ 与 $\Delta^\circ$ 如下计算：
\begin{equation}\label{eq:toric-hodge-number}
h_{\mathrm{toric}}^{1,1}(V)
=l(\Delta^\circ)-n-1-\sum_{\Gamma^\circ}l^*(\Gamma^\circ),
\end{equation}
\begin{equation}\label{eq:polynomial-hodge-number}
h_{\mathrm{poly}}^{n-2,1}(V)
=l(\Delta)-n-1-\sum_\Gamma l^*(\Gamma),
\end{equation}
其中求和分别遍历 $\Delta^\circ$ 与 $\Delta$ 的所有刻面 $\Gamma^\circ$ 和 $\Gamma$；$l$、$l^*$ 如式~\eqref{eq:lattice-points} 定义。一旦证明这个断言，便立即得到等式
\[
h_{\mathrm{toric}}^{1,1}(V)=h_{\mathrm{poly}}^{n-2,1}(V^\circ),
\qquad
h_{\mathrm{toric}}^{1,1}(V^\circ)=h_{\mathrm{poly}}^{n-2,1}(V),
\]
因为由 $V$ 换到 $V^\circ$，正好交换了式~\eqref{eq:toric-hodge-number}、\eqref{eq:polynomial-hodge-number} 中 $\Delta$ 与 $\Delta^\circ$ 的角色。还应提到，可以构造自然同构 $H_{\mathrm{toric}}^{1,1}(V)\simeq H_{\mathrm{poly}}^{n-2,1}(V^\circ)$，称为\emph{单项式--\IndexLocation{indexloc:0910:80}除子镜像映射}，见 \SourceCite{AGM1}。

先证明式~\eqref{eq:toric-hodge-number}。设 $v_i\in\Sigma(1)$ 是 $\Delta^\circ$ 的某个刻面的内部格点。环面爆破的标准描述表明，$f(D_i)$ 是一个点，其中 $f:X_\Sigma\to\mathbb P_\Delta$ 为自然映射。立即可知，在 $X_\Sigma$ 中，$D_i$ 与一般的 $V\in|-K_{X_\Sigma}|$ 不相交。因此这些元素对 $H_{\mathrm{toric}}^{1,1}(V)$ 没有贡献。置
\[
\Sigma(1)'=\{v_i\in\Sigma(1):
v_i\text{ 不在 }\Delta^\circ\text{ 的任何刻面的相对内部}\},
\]
则 $H_{\mathrm{toric}}^{1,1}(V)$ 由 $v_i\in\Sigma(1)'$ 对应的 $D_i'$ 生成。要找出各 $D_i'$ 之间的关系，注意 $m\in M$ 给出 $T_N$ 的特征标 $\chi^m$，从而给出
\SourcePage{81}
$X_\Sigma$ 上的有理函数。式~\eqref{eq:toric-principal-divisor} 已给出 $\chi^m$ 的除子；将该公式限制到 $V$，得到线性等价关系
\begin{equation}\label{eq:toric-restricted-divisor-relation}
\sum_{v_i\in\Sigma(1)'}\langle m,v_i\rangle D_i'\sim0.
\end{equation}
\SourceCite{AGM1} 证明，这些就是各 $D_i'$ 之间的全部关系。因此
\[
H_{\mathrm{toric}}^{1,1}(V)
\simeq\bigl(\mathbb Z^{\Sigma(1)'}/M\bigr)\otimes\mathbb C,
\]
其中嵌入 $M\to\mathbb Z^{\Sigma(1)'}$ 按式~\eqref{eq:toric-restricted-divisor-relation} 的要求，由 $m\mapsto(\langle m,v_i\rangle)$ 给出。于是
\[
h_{\mathrm{toric}}^{1,1}(V)=|\Sigma(1)'|-n.
\]
用式~\eqref{eq:lattice-points} 的记号，因为 $\Sigma(1)=\Delta^\circ\cap N-\{0\}$，可见 $|\Sigma(1)|=l(\Delta^\circ)-1$；这里用到了射影细分 $\Sigma$ 的极大性。去掉各刻面的内部格点后，得到
$|\Sigma(1)'|=l(\Delta^\circ)-1-\sum_{\Gamma^\circ}l^*(\Gamma^\circ)$，式~\eqref{eq:toric-hodge-number} 随即成立。

现在考虑 $h_{\mathrm{poly}}^{n-2,1}(V)$。我们将描述来自族 $V\in|-K_{X_\Sigma}|$ 的模。这里先作非正式的讨论，\hyperref[chap:moduli]{第~\ref*{chap:moduli} 章}会作仔细处理。引理~\ref{lem:projective-subdivision-properties} 表明，$\Delta$ 是 $-K_{X_\Sigma}$ 对应的多胞形。因此，Laurent 多项式 $f\in L(\Delta)$ 给出仿射超曲面 $Z_f\subset T_N$，而且可设 $V=\overline Z_f\subset X_\Sigma$。由于将 $f$ 乘以非零常数不改变 $V$，变动 $f$ 得到一个维数为 $l(\Delta)-1$ 的多项式形变参数空间。不过，必须计入 $X_\Sigma$ 的自同构。因为 $\Delta$ 对应于 $-K_{X_\Sigma}$，命题~\ref{prop:automorphism-facets} 表明，自同构群的维数为
\[
\dim\operatorname{Aut}(X_\Sigma)=n+\sum_\Gamma l^*(\Gamma).
\]
现在容易推出，多项式模的维数由式~\eqref{eq:polynomial-hodge-number} 给出。

要完成定理的证明，还需理解 $V$ 的“非环面”除子类与“非多项式”形变。Batyrev \SourceCite{Batyrev4} 证明，在这两种情形中，“修正项”均由 $\Delta$ 与 $\Delta^\circ$ 的余维数为二的面确定。公式为
\begin{equation}\label{eq:batyrev-h11}
\begin{aligned}
h^{1,1}(V)={}&l(\Delta^\circ)-n-1
-\sum_{\Gamma^\circ}l^*(\Gamma^\circ)\\
&+\sum_{\Theta^\circ}l^*(\Theta^\circ)l^*(\widehat\Theta^\circ),
\end{aligned}
\end{equation}
\begin{equation}\label{eq:batyrev-hn21}
\begin{aligned}
h^{n-2,1}(V)={}&l(\Delta)-n-1-\sum_\Gamma l^*(\Gamma)\\
&+\sum_\Theta l^*(\Theta)l^*(\widehat\Theta),
\end{aligned}
\end{equation}
其中第一个等式里的 $\Gamma^\circ$ 指 $\Delta^\circ$ 的余维数为一的面，$\Theta^\circ$ 指其余维数为二的面，$\widehat\Theta^\circ$ 指 $\Delta$ 中与 $\Theta^\circ$ 对偶的面。第二个等式的含义类似。如前所述，将这些公式分别应用于 $V$ 和 $V^\circ$，就交换了 $\Delta$ 与 $\Delta^\circ$ 的角色，因此定理立即由式~\eqref{eq:batyrev-h11}、\eqref{eq:batyrev-hn21} 推出。

式~\eqref{eq:batyrev-h11} 的修正项来自与一般 $V$ 的交可约的除子 $D_i$。前面已看到，可以忽略那些 $v_i$ 位于刻面 $\Gamma^\circ$ 相对内部的 $D_i$。一般地，设 $v_i$ 位于某个面 $\Theta^\circ\subset\Delta^\circ$ 的相对内部。若 $\widehat\Theta^\circ\subset\Delta$ 表示对偶面，则 $\mathbb P_{\widehat\Theta^\circ}\subset\mathbb P_\Delta$ 是轨道闭包，并且在爆破映射 $f:X_\Sigma\to\mathbb P_\Delta$ 下，有 $f(D_i)=\mathbb P_{\widehat\Theta^\circ}$。若 $\overline V=f(V)$ 是 $\mathbb P_\Delta$ 中相应的反典范超曲面，则 $D_i'=V\cap D_i$ 与 $\overline V\cap\mathbb P_{\widehat\Theta^\circ}$ 有相同的不可约分支数。

回顾，面 $\Theta^\circ$ 与其对偶 $\widehat\Theta^\circ$ 满足 $\dim\Theta^\circ+\dim\widehat\Theta^\circ=n-1$。因此，若 $\Theta^\circ$ 的余维数至少为三，则 $\widehat\Theta^\circ$ 的维数至少为二，所以 $\mathbb P_{\widehat\Theta^\circ}$ 的
\SourcePage{82}
维数至少为二。Bertini 定理表明 $\overline V\cap\mathbb P_{\widehat\Theta^\circ}$ 不可约。因此，当 $v_i$ 位于 $\Delta^\circ$ 的余维数至少为三的面中时，$D_i'$ 不可约。

另一方面，若 $\Theta^\circ$ 的余维数为二，则 $\widehat\Theta^\circ$ 的维数为一，$\mathbb P_{\widehat\Theta^\circ}$ 是一条曲线。在这种情形，\SourceCite[第 5.3 节]{Fulton3} 表明
\[
\overline V\cdot\mathbb P_{\widehat\Theta^\circ}
=\operatorname{vol}(\widehat\Theta^\circ)
=l^*(\widehat\Theta^\circ)+1,
\]
这里 $\operatorname{vol}$ 表示归一化体积；第二个等式成立，是因为 $\widehat\Theta^\circ$ 为一维。这说明 $\overline V\cap\mathbb P_{\widehat\Theta^\circ}$ 由 $l^*(\widehat\Theta^\circ)+1$ 个点组成，从而 $D_i'$ 有 $l^*(\widehat\Theta^\circ)+1$ 个连通分支。可以证明，这些分支在 $H^{1,1}(V)$ 中给出线性无关的类。它们的和 $D_i'$ 已经属于 $H_{\mathrm{toric}}^{1,1}(V)$，因此，$\Theta^\circ$ 的每个内部格点增加 $l^*(\widehat\Theta^\circ)$ 个新的类。因为 $\Sigma$ 极大，使用了 $\Theta^\circ$ 的全部内部格点，这就解释了式~\eqref{eq:batyrev-h11} 中的修正项。

至于式~\eqref{eq:batyrev-hn21} 的修正项，我们将解释这个公式为什么合理；完整证明参见 \SourceCite{Batyrev4,Voisin3}。令 $Z_f=V\cap T_N\subset T_N$ 为 $V$ 所确定的仿射超曲面，则补集 $V-Z_f=\sum_i V\cap D_i$ 是轨形意义下的正规交叉除子。因此，Gysin 谱序列给出正合列
\[
\begin{aligned}
\longrightarrow\bigoplus_i H^{n-3}(V\cap D_i,\mathbb C)(-1)
&\longrightarrow H^{n-1}(V,\mathbb C)\\
&\longrightarrow\operatorname{Gr}_{n-1}^W H^{n-1}(Z_f,\mathbb C)
\longrightarrow0.
\end{aligned}
\]
这些群都带有纯 Hodge 结构；由此可以证明
\begin{equation}\label{eq:gysin-hodge-number}
\begin{aligned}
h^{n-2,1}(V)={}&h^{n-2,1}
\bigl(\operatorname{Gr}_{n-1}^W H^{n-1}(Z_f,\mathbb C)\bigr)\\
&+\sum_i h^{n-3,0}(V\cap D_i).
\end{aligned}
\end{equation}
根据 Danilov 与 Khovanskii 的工作 \SourceCite[5.9]{DK}，右边第一项可计算为
\[
\begin{aligned}
&h^{n-2,1}\bigl(\operatorname{Gr}_{n-1}^W H^{n-1}(Z_f,\mathbb C)\bigr)\\
&\qquad=l^*(2\Delta)-(n+1)l^*(\Delta)-\sum_\Gamma l^*(\Gamma).
\end{aligned}
\]
不过，自反性意味着 $l^*(\Delta)=1$，而 $l^*(2\Delta)=l(\Delta)$ 也来自自反性。因此上式简化为
\[
\begin{aligned}
&h^{n-2,1}\bigl(\operatorname{Gr}_{n-1}^W H^{n-1}(Z_f,\mathbb C)\bigr)\\
&\qquad=l(\Delta)-n-1-\sum_\Gamma l^*(\Gamma).
\end{aligned}
\]
余下只须证明，式~\eqref{eq:gysin-hodge-number} 右边第二项等于式~\eqref{eq:batyrev-hn21} 中的修正项。

理解它的一种方式如下。设 $v_i$ 位于面 $\widehat\Theta\subset\Delta^\circ$ 的相对内部，令 $\Theta\subset\Delta$ 表示对偶面。如上，这给出轨道闭包 $\mathbb P_\Theta\subset\mathbb P_\Delta$，且 $f(D_i)=\mathbb P_\Theta$。于是有映射 $V\cap D_i\to\overline V\cap\mathbb P_\Theta$，进一步诱导自然映射
\[
H^{n-3,0}(\overline V\cap\mathbb P_\Theta)
\longrightarrow H^{n-3,0}(V\cap D_i).
\]
大多数情形下，$h^{n-3,0}(\overline V\cap\mathbb P_\Theta)=0$；但当 $\Theta$ 单纯，且 $\overline V\cap\mathbb P_\Theta$ 的维数为 $n-3$ 时，\SourceCite[5.5]{DK} 给出公式
\begin{equation}\label{eq:face-holomorphic-forms}
h^{n-3,0}(\overline V\cap\mathbb P_\Theta)=l^*(\Theta).
\end{equation}
注意，$\dim(\overline V\cap\mathbb P_\Theta)=n-3$ 等价于 $\Theta$ 是余维数为二的面。因此，若 $\Theta\subset\Delta$ 是单纯的余维数为二的面，则 $\widehat\Theta$ 的 $l^*(\widehat\Theta)$ 个内部格点 $v_i$ 中的每一个，都贡献 $l^*(\Theta)$ 个全纯 $(n-3)$ 形式。这样，修正项 $\sum_\Theta l^*(\Theta)l^*(\widehat\Theta)$ 就自然出现了。遗憾的是，这个论证并不适用于一般情形。为实际计算式~\eqref{eq:gysin-hodge-number} 右边的和，需要将 $D_i$ 分解为轨道 $O$ 的不交并，并用
\SourcePage{83}
\SourceCite{DK} 的方法计算各 $V\cap O$ 的适当 $h^{p,q}$ 数。如上所述，这一论证的细节可在 \SourceCite{Batyrev4,Voisin3} 中找到。
\end{proof}

从历史上看，人们最初将大量 Calabi--Yau 簇族构造成加权射影空间中的拟光滑超曲面；最终找到了 $7555$ 个例子 \SourceCite{KS,KS2}。将这些例子的 Hodge 数绘成图，所得图形显示出醒目的\emph{部分对称性}，这是最早支持镜像对称的一些证据。\SourceCite{Batyrev4} 发表后，\SourceCite{CdK} 核查了表中全部 $7555$ 个簇：一般 Calabi--Yau 超曲面的 Newton 多胞形，经平移使唯一内部格点位于原点后，的确都是自反的。因此这些例子都是 Batyrev 镜像构造的特殊情形。\hyperref[subsec:nonreflexive-weights]{第 4.4.3 小节}将进一步介绍所用的方法。

定理~\ref{thm:batyrev-hodge} 给出的 Hodge 数相等关系，例如说明 $H^{1,1}(V)$ 同构于 $H^1(V^\circ,T_{V^\circ})\simeq H^{n-2,1}(V^\circ)$，但并未指定同构。根据镜像对称，联系 $V$ 的 Kähler 模与 $V^\circ$ 的复模的镜像映射，应至少定义在 $V$ 的 Kähler 模空间的一个开子集上；其导数应给出典范同构 $H^{1,1}(V)\simeq H^{n-2,1}(V^\circ)$。此外，在上述证明中还提到了单项式--除子镜像映射，即自然同构 $H_{\mathrm{toric}}^{1,1}(V)\simeq H_{\mathrm{poly}}^{n-2,1}(V^\circ)$。\hyperref[chap:moduli]{第~\ref*{chap:moduli} 章}将更详细地讨论此同构，并解释它与镜像映射的关系。\index{单项式--除子镜像映射|hyperpage}\index{Hodge数@Hodge 数|hyperpage}

还应提到，根据式~\eqref{eq:mirror-cohomology-forms}，镜像对称暗示更一般地，当 $V^\circ$ 是 $V$ 的 Batyrev 对偶时，应有 Hodge 数相等关系
\begin{equation}\label{eq:full-mirror-hodge}
h^{p,q}(V)=h^{n-1-p,q}(V^\circ).
\end{equation}
回顾，此时 $V$ 与 $V^\circ$ 的维数为 $d=n-1$。定理~\ref{thm:batyrev-hodge} 包含了 $(p,q)=(1,1)$ 与 $(n-2,1)$ 的情形，而 $q=0$ 的情形是平凡的。对三维簇，这些就是所需的全部情形；但在维数至少为四时，式~\eqref{eq:full-mirror-hodge} 是否成立仍是一个未解决的问题。

当 $V\subset X_\Sigma$ 是我们正在考虑的这类环面 Calabi--Yau 超曲面时，\SourceCite{BB2,BD} 定义了\emph{弦论 Hodge 数} $h_{\mathrm{st}}^{p,q}(V)$。当 $V$ 光滑或 $q=1$ 时，这些数与通常的 Hodge 数 $h^{p,q}(V)$ 相同。更重要的是，当 $V^\circ$ 是 $V$ 的 Batyrev 对偶时，对所有 $p+q=\dim V$，这些数满足
\[
h_{\mathrm{st}}^{p,q}(V)=h_{\mathrm{st}}^{n-1-p,q}(V^\circ).
\]
因此，在光滑情形且 $p+q=\dim V$ 时，式~\eqref{eq:full-mirror-hodge} 成立；不过，如前所述，若 $V$ 或 $V^\circ$ 是有奇点的极小 Calabi--Yau 簇，这个问题仍未解决。\index{弦论Hodge数@弦论 Hodge 数|hyperpage}

\subsection{多重镜像}\label{subsec:multiple-mirrors}

在 Batyrev 镜像构造中，$\Sigma$ 与 $\Sigma^\circ$ 的极大射影细分可能有不止一种选择。这似乎导致多个镜像族。例如，若 $\Sigma$ 与 $\Sigma'$ 都满足 $\Sigma(1)=\Sigma'(1)=\Delta^\circ\cap N-\{0\}$，便得到两个族 $V\subset X_\Sigma$ 与 $V'\subset X_{\Sigma'}$。$V^\circ\subset X_{\Sigma^\circ}$ 能否同时为这两个族的镜像？答案是肯定的，这是 \SourceCite{AGM2} 讨论的“多重镜像”现象的一部分。\hyperref[chap:moduli]{第~\ref*{chap:moduli} 章}将详细讨论多重镜像。不过，若限于式~\eqref{eq:toric-hodge-subspace} 中的环面 Kähler 模，大致的想法是：$\Sigma$ 与 $\Sigma'$ 对应于 GKZ 分解中不同的锥，而族 $V$ 与 $V'$ 的 Kähler 模给出\emph{同一个}整体 Kähler 模空间中互不相交的开集。\index{多重镜像|hyperpage}

\SourcePage{84}
也可以在另一侧作不同选择。此时，极大射影细分 $\Sigma^\circ$ 与 $\Sigma'^\circ$ 给出族 $V^\circ$ 与 $V'^\circ$。我们希望这些族的复模对应于族 $V$ 的 Kähler 模。因此 $V^\circ$ 与 $V'^\circ$ 应有\emph{相同的}复模。若限于式~\eqref{eq:polynomial-hodge-subspace} 中的多项式形变，这一点来自定理~\ref{thm:batyrev-hodge} 的证明。具体说，$V^\circ$ 的多项式模对应于 $L(\Delta^\circ)$ 的一个开子集模去 $X_{\Sigma^\circ}$ 的自同构群的作用。不过，$V'^\circ$ 也由同一空间 $L(\Delta^\circ)$ 中的多项式定义；利用 \SourceCite{Cox} 的结果，可以看到 $\operatorname{Aut}(X_{\Sigma^\circ})$ 与 $\operatorname{Aut}(X_{\Sigma'^\circ})$ 有同样的单位元连通分支。因此，至少在局部，这些多项式模空间看起来相同。\hyperref[chap:moduli]{第~\ref*{chap:moduli} 章}将仔细解释这些内容。

\section{重访五次三维簇}\label{sec:quintic-revisited}

现在回到\hyperref[chap:quintic]{第 2 章}讨论的五次三维簇 $V\subset\mathbb P^4$。由于 $V$ 是 Fano 环面簇 $\mathbb P^4$ 中的反典范超曲面，可用 Batyrev 构造描述它的镜像族。

$\mathbb P^4$ 上的反典范类 $D_0+\cdots+D_4$ 对应于\IndexLocation{indexloc:1177:84}\IndexLocation{indexloc:1223:84}自反多胞形 $\Delta\subset M_{\mathbb R}\simeq\mathbb R^4$，其顶点为
\begin{equation}\label{eq:quintic-polytope-vertices}
\begin{gathered}
(-1,-1,-1,-1),\\
(4,-1,-1,-1),\\
(-1,4,-1,-1),\\
(-1,-1,4,-1),\\
(-1,-1,-1,4).
\end{gathered}
\end{equation}
可以把 $\Delta$ 视为 $\mathbb P^4$ 中五次多项式的 Newton 多胞形。为看出这一点，考虑如式~\eqref{eq:laurent-space} 那样由 $\Delta$ 确定的 Laurent 多项式，例如
\[
f=\sum_{i=1}^s a_it^{m_i},\qquad
\Delta\cap M=\{m_1,\ldots,m_s\}.
\]
再注意，由式~\eqref{eq:toric-graded-sections} 给出的齐次化
\[
F(x_0,x_1,x_2,x_3,x_4)=
x_0x_1x_2x_3x_4\,f(x_1/x_0,x_2/x_0,x_3/x_0,x_4/x_0)
\]
就是 $x_0,x_1,x_2,x_3,x_4$ 的任意五次多项式。

为寻找 $V$ 的镜像族，使用对偶多胞形 $\Delta^\circ\subset N_{\mathbb R}$。容易看到，它的顶点为
\begin{equation}\label{eq:quintic-polar-vertices}
\begin{gathered}
(-1,-1,-1,-1),\\
(1,0,0,0),\\
(0,1,0,0),\\
(0,0,1,0),\\
(0,0,0,1).
\end{gathered}
\end{equation}
环面簇 $\mathbb P_{\Delta^\circ}$ 由 $\Delta^\circ$ 在 $M_{\mathbb R}$ 中的法扇确定，该扇由向 $\Delta$ 的各面作出的锥组成。因此，一维锥生成元是式~\eqref{eq:quintic-polytope-vertices} 中的向量，它们生成 $M$ 中指数为 $125$ 的子格 $M_1$。商 $M/M_1$ 为
\[
G=\left\{(a_0,a_1,a_2,a_3,a_4)\in(\mathbb Z_5)^5:
\sum_{i=0}^4a_i=0\pmod5\right\}/\mathbb Z_5,
\]
其中 $\mathbb Z_5\subset(\mathbb Z_5)^5$ 是对角子群。由于 $M_{\mathbb R}=(M_1)_{\mathbb R}$，可将 $\Delta^\circ$ 的法扇视为 $(M_1)_{\mathbb R}$ 中的扇。其生成元由 $M_1$ 的一组 $\mathbb Z$ 基以及这些基向量之和的相反数组成。这说明，使用格 $M_1$ 时，该法扇就是 $\mathbb P^4$ 的标准扇。若再换到更大的
\SourcePage{85}
格 $M$，环面几何的标准结果表明，$\mathbb P_{\Delta^\circ}$ 是 $\mathbb P^4$ 模去 $M/M_1\simeq G$ 的商，其中 $G$ 在 $\mathbb P^4$ 上的作用与\hyperref[chap:quintic]{第 2 章}所述完全相同。

$\mathbb P_{\Delta^\circ}$ 的齐次坐标环 $S$ 可描述如下。它是多项式环 $S=\mathbb C[x_0,x_1,x_2,x_3,x_4]$，由 Chow 群 $A_3(\mathbb P_{\Delta^\circ})$ 分次。由式~\eqref{eq:toric-chow-sequence} 得到正合列
\[
0\longrightarrow\mathbb Z^4\longrightarrow\mathbb Z^5
\longrightarrow\mathbb Z\oplus G\longrightarrow0,
\]
其中第一个映射的矩阵以式~\eqref{eq:quintic-polytope-vertices} 的各行为行，第二个映射将 $(a_0,a_1,a_2,a_3,a_4)\in\mathbb Z^5$ 送到
\[
\begin{aligned}
\bigl[a_0,a_1,a_2,a_3,a_4\bigr]
={}&\left(\sum_{i=0}^4a_i,
(-a_1-a_2-a_3-a_4,a_1,a_2,a_3,a_4)\right)\!\in\mathbb Z\oplus G.
\end{aligned}
\]
因此 $A_3(\mathbb P_{\Delta^\circ})\simeq\mathbb Z\oplus G$，而 $S$ 的分次规定单项式 $x_0^{a_0}\cdots x_4^{a_4}$ 的次数为 $[a_0,a_1,a_2,a_3,a_4]\in\mathbb Z\oplus G$。

于是，$\mathbb P_{\Delta^\circ}$ 的反典范类的次数为 $[1,1,1,1,1]=(5,0)\in\mathbb Z\oplus G$，$S$ 中仅有的这个次数的单项式是 $x_i^5$ 及 $x_0x_1x_2x_3x_4$。因此得到超曲面 $\overline V^\circ\subset\mathbb P_{\Delta^\circ}$，由方程
\begin{equation}\label{eq:quintic-general-mirror}
a_0x_0^5+a_1x_1^5+a_2x_2^5+a_3x_3^5+a_4x_4^5
+a_5x_0x_1x_2x_3x_4=0
\end{equation}
定义。具体说，$\overline V^\circ$ 就是 $\mathbb P^4$ 中由上述方程定义的超曲面模去 $G$ 的商。注意，式~\eqref{eq:quintic-general-mirror} 中的单项式对应于 $\Delta^\circ$ 中的整点，即原点与式~\eqref{eq:quintic-polar-vertices} 中的各点。此外，关系
\[
(x_0^5)(x_1^5)(x_2^5)(x_3^5)(x_4^5)
=(x_0x_1x_2x_3x_4)^5
\]
反映了式~\eqref{eq:quintic-polar-vertices} 中五个向量的和等于原点的五倍。

方程~\eqref{eq:quintic-general-mirror} 不同于\hyperref[chap:quintic]{第 2 章}讨论的单参数簇族。不过，回顾我们关心的是这些超曲面的模。$\mathbb P_{\Delta^\circ}$ 的\IndexLocation{indexloc:1355:85}自同构群是环面 $T_M$；利用这些自同构，可以约化到\hyperref[chap:quintic]{第 2 章}考虑的单参数族
\[
x_0^5+x_1^5+x_2^5+x_3^5+x_4^5+\psi x_0x_1x_2x_3x_4=0.
\]
\hyperref[chap:moduli]{第~\ref*{chap:moduli} 章}将给出完成此约化的系统方法。

最后，为得到五次超曲面 $V$ 的镜像，在 $M_{\mathbb R}$ 中选取\IndexLocation{indexloc:1096:85}极大射影细分 $\Sigma^\circ$。这样的扇细化 $\Delta^\circ$ 的法扇，并满足 $\Sigma^\circ(1)=\Delta\cap M-\{0\}$。于是镜像族由超曲面 $V^\circ\subset X_{\Sigma^\circ}$ 给出，它们是 $\overline V^\circ\subset\mathbb P_{\Delta^\circ}$ 的严格变换。

注意，$M\cap\Delta-\{0\}$ 除顶点外还有大量点，因此 $\Sigma^\circ$ 有许多选择。找出所有这样的 $\Sigma^\circ$ 这一计算问题，远远超出了当时计算机的能力。\SourceCite{Morrison2} 的附录描述了 $\Sigma^\circ$ 的一个具体选择。

\section{环面完全交}\label{sec:toric-complete-intersections}

推广\hyperref[sec:batyrev-mirror]{第 4.1 节}研究的环面超曲面，也可以问环面簇中的 Calabi--Yau 完全交是否具有镜像对称。本节通过介绍 Batyrev 与 Borisov \SourceCite{Borisov1,BB1} 的两种镜像对称构造来回答这个问题；这些构造涉及 nef 分拆与自反 Gorenstein 锥。

\SourcePage{86}
\subsection{nef 分拆}\label{subsec:nef-partitions}

设 $X=\mathbb P_\Delta$ 是自反多胞形 $\Delta$ 对应的 $n$ 维环面簇。令 $D_\rho$ 表示 $\rho\in\Sigma(1)$ 对应的环面除子，其中 $\Sigma$ 是 $\Delta$ 的法扇。给定将 $\Sigma(1)=I_1\cup\cdots\cup I_k$ 分成 $k$ 个互不相交的子集的分拆，得到除子 $E_j=\sum_{\rho\in I_j}D_\rho$，满足 $-K_X=E_1+\cdots+E_k$。
\begin{definition}\label{def:nef-partition}
若对每个 $j$，$E_j$ 都是由整体截面生成的 Cartier 除子，则称分解 $\Sigma(1)=I_1\cup\cdots\cup I_k$ 为 \emph{nef 分拆}。
\end{definition}

若 $\Delta_j$ 是 $E_j$ 对应的多胞形，则 nef 分拆意味着
\[
\Delta=\Delta_1+\cdots+\Delta_k.
\]
不严格地说，分解 $-K_X=E_1+\cdots+E_k$ 与 $\Delta=\Delta_1+\cdots+\Delta_k$ 也称为 nef 分拆。

$\mathcal O(E_1),\ldots,\mathcal O(E_k)$ 的一般截面定义一个维数为 $d=n-k$ 的完全交 $\overline V\subset X$。利用 $-K_X=E_1+\cdots+E_k$ 与伴随公式，看到 $\overline V\subset\mathbb P_\Delta$ 是 Calabi--Yau 簇。如\hyperref[sec:batyrev-mirror]{第 4.1 节}，$\Delta$ 的法扇的极大射影细分 $\Sigma'$ 给出映射 $X_{\Sigma'}\to\mathbb P_\Delta$，而 $\overline V\subset\mathbb P_\Delta$ 的严格变换 $V\subset X_{\Sigma'}$ 仍是完全交，并且成为极小 Calabi--Yau 簇，见 \SourceCite[命题 4.15]{BB1}。

为描述 $V$ 的镜像族，仍将 $\Sigma(1)$ 中的元素视为一维锥生成元。在自反情形，它们恰好是极多胞形 $\Delta^\circ\subset N_{\mathbb R}$ 的顶点。于是 nef 分拆 $\Sigma(1)=I_1\cup\cdots\cup I_k$ 给出多胞形
\[
\nabla_j=\operatorname{Conv}(\{0\}\cup I_j)\subset N_{\mathbb R}.
\]
立即得到
\[
\Delta^\circ=\operatorname{Conv}(\nabla_1\cup\cdots\cup\nabla_k).
\]
\SourceCite{Borisov1} 的基本结果是：Minkowski 和
\[
\nabla=\nabla_1+\cdots+\nabla_k
\]
仍是自反多胞形，而且 $\nabla=\nabla_1+\cdots+\nabla_k$ 也是 nef 分拆。按 \SourceCite{BB1} 的术语，这称为\emph{对偶 nef 分拆}。

使用对偶 nef 分拆，得到 Calabi--Yau 完全交 $\overline V^\circ\subset\mathbb P_\nabla$；选取 $\nabla$ 的法扇的极大射影细分 $\Sigma^\circ$，得到极小 Calabi--Yau 完全交族 $V^\circ\subset X_{\Sigma^\circ}$。猜想为
\[
V^\circ\text{ 是 }V\text{ 的镜像族。}
\]
这个猜想与 \SourceCite{LTe} 提出的射影空间完全交的镜像对称构造一致。

作为支持此猜想的证据，Batyrev 与 Borisov 证明
\begin{equation}\label{eq:complete-intersection-hodge}
h^{p,1}(V)=h^{d-p,1}(V^\circ),
\end{equation}
其中 $d=n-k$ 是 $V$ 与 $V^\circ$ 的维数。更一般地，使用\hyperref[subsec:batyrev-hodge]{第 4.1.3 小节}讨论的弦论 Hodge 数 $h_{\mathrm{st}}^{p,q}$，可以证明
\begin{equation}\label{eq:complete-intersection-string-hodge}
h_{\mathrm{st}}^{p,q}(V)=h_{\mathrm{st}}^{d-p,q}(V^\circ),
\end{equation}
见 \SourceCite{BB2}。由于 $h^{p,1}(V)=h_{\mathrm{st}}^{p,1}(V)$，可见式~\eqref{eq:complete-intersection-hodge} 是式~\eqref{eq:complete-intersection-string-hodge} 的特殊情形。\index{nef分拆@nef 分拆|hyperpage}\index{完全交|hyperpage}

\SourcePage{87}
\subsection{Gorenstein 锥}\label{subsec:gorenstein-cones}

研究 Gorenstein 锥，可推广 nef 分拆的组合对偶性。其定义如下。设有格 $\overline M$ 及其对偶 $\overline N$。
\begin{definition}\label{def:gorenstein-cone}
设 $\sigma\subset\overline M_{\mathbb R}$ 是有理多面锥。
\begin{enumerate}[label=(\roman*)]
\item 若 $\sigma$ 强凸，且存在 $n_\sigma\in\overline N$，使 $\sigma$ 的每个本原整生成元 $v$ 都满足 $\langle v,n_\sigma\rangle=1$，则称 $\sigma$ 为 \emph{Gorenstein 锥}。
\item 若 $\sigma$ 及其对偶 $\check\sigma$ 都是 Gorenstein 锥，则称 $\sigma$ 为\emph{自反 Gorenstein 锥}。
\end{enumerate}
\end{definition}
显然，$\sigma$ 是自反 Gorenstein 锥，当且仅当 $\check\sigma$ 是。此外，若 $\sigma$ 自反，则它必有满维数，这意味着 $n_\sigma\in\overline N$ 唯一确定。$\check\sigma$ 相应地确定 $m_\sigma\in\overline M$，整数
\[
r=\langle m_\sigma,n_\sigma\rangle
\]
称为 $\sigma$，或 $\check\sigma$，的\emph{指数}。

先解释自反 Gorenstein 锥与自反多胞形对偶性的关系。给定以原点为内点的整多胞形 $\Delta\subset M_{\mathbb R}$，将格 $M$ 扩充为 $\overline M=M\times\mathbb Z$，相应地有 $\overline N=N\times\mathbb Z$。令 $\sigma\subset\overline M_{\mathbb R}$ 为由元素 $\{(v,1):v\in\Delta\}$ 张成的锥。换言之，
\begin{equation}\label{eq:polytope-gorenstein-cone}
\sigma=\{(\lambda v,\lambda):v\in\Delta,\ \lambda\geq0\}.
\end{equation}
注意，$\sigma$ 是 Gorenstein 锥。简单论证表明，$\sigma$ 自反，当且仅当 $\Delta$ 自反。此外，$\Delta$ 自反时，指数 $r=1$，并且 $\check\sigma$ 由 $\Delta^\circ$ 构造的方式，与 $\sigma$ 由 $\Delta$ 构造的方式相同。

也可用自反 Gorenstein 锥描述 nef 分拆的对偶性。给定自反多胞形 $\Delta\subset M_{\mathbb R}$ 及 nef 分拆 $\Delta=\Delta_1+\cdots+\Delta_k$，定义格 $\overline M=M\times\mathbb Z^k$，其对偶格为 $\overline N=N\times\mathbb Z^k$。令 $\sigma$ 为锥
\begin{equation}\label{eq:nef-gorenstein-cone}
\sigma=\left\{\left(\sum_{i=1}^k\lambda_iv_i,
\lambda_1,\ldots,\lambda_k\right):
v_i\in\Delta_i,\ \lambda_i\geq0\right\}.
\end{equation}
注意，它推广了式~\eqref{eq:polytope-gorenstein-cone}。

\begin{proposition}\label{prop:nef-cone-duality}
$\sigma$ 是指数为 $k$ 的自反 Gorenstein 锥。它的对偶锥 $\check\sigma$，由对偶 nef 分拆 $\nabla=\nabla_1+\cdots+\nabla_k$ 应用式~\eqref{eq:nef-gorenstein-cone} 构造得到。
\end{proposition}
\begin{proof}
证明见 \SourceCite{BB1}。
\end{proof}

当 $\sigma$ 如式~\eqref{eq:nef-gorenstein-cone} 定义时，对偶锥 $\check\sigma$ 具有有趣的几何意义。这个锥是某个扇 $\overline\Sigma$ 的支集，其对应环面簇为 $X=\mathbb P_\Delta$ 上向量丛
\begin{equation}\label{eq:nef-vector-bundle}
\mathcal O(-E_1)\oplus\cdots\oplus\mathcal O(-E_k)
\end{equation}
的总空间。环面簇 $X$ 可视为经零截面嵌入 $X_{\overline\Sigma}$，因此 $V$ 可以构造成 $X_{\overline\Sigma}$ 的子簇。$\overline\Sigma$ 也可扩充为完备扇，其对应环面簇是 $X$ 上的射影丛
\[
\mathbb P\bigl(\mathcal O\oplus\mathcal O(-E_1)\oplus\cdots\oplus\mathcal O(-E_k)\bigr).
\]
给出完全交 $V$ 的 $\mathcal O(E_1),\ldots,\mathcal O(E_k)$ 的各截面，也确定了向量丛 $\mathcal O(E_1)\oplus\cdots\oplus\mathcal O(E_k)$ 的一个整体截面。相应的超曲面为
\begin{equation}\label{eq:cayley-hypersurface}
Y'\subset\mathbb P\bigl(\mathcal O(E_1)\oplus\cdots\oplus\mathcal O(E_k)\bigr)
\end{equation}
这个超曲面与 $V$ 密切相关。这就是 \SourceCite{BB3,GKZ2} 描述的 \emph{Cayley 技巧}。\index{Gorenstein锥@Gorenstein 锥|hyperpage}\index{Cayley技巧@Cayley 技巧|hyperpage}

\SourcePage{88}
式~\eqref{eq:cayley-hypersurface} 的超曲面 $Y'$ 不是 Calabi--Yau 簇，但它是 Schimmrigk 意义下的\emph{广义 Calabi--Yau 簇}，见 \SourceCite{BB3,Schimmrigk}。可用任意自反 Gorenstein 锥 $\sigma$ 如下构造广义 Calabi--Yau 簇。由 $\sigma$ 得到半群 $\sigma\cap\overline M$，并将 $m\in\sigma\cap\overline M$ 的次数定义为 $\langle m,n_\sigma\rangle$。这使 $\mathbb C[\sigma\cap\overline M]$ 成为分次环，给出环面簇
\[
\mathbb P_\sigma=\operatorname{Proj}\bigl(\mathbb C[\sigma\cap\overline M]\bigr).
\]
若 $\mathcal O_{\mathbb P_\sigma}(1)$ 为 $\mathbb P_\sigma$ 上的重言层，则它的一个整体截面定义子簇
\[
Y\subset\mathbb P_\sigma,
\]
Batyrev 与 Borisov 称其为\emph{与 $\sigma$ 对应的广义 Calabi--Yau 簇}。当然，对偶锥 $\check\sigma$ 也是自反 Gorenstein 锥，因此得到对偶族
\[
Y^\circ\subset\mathbb P_{\check\sigma}.
\]
猜想它在适当意义下是 $Y\subset\mathbb P_\sigma$ 的镜像，见 \SourceCite[猜想 2.17]{BB3}。

这个构造包含了目前考虑过的 Calabi--Yau 超曲面与完全交。更精确地说：
\begin{itemize}
\item 当 $\sigma$ 由式~\eqref{eq:polytope-gorenstein-cone} 给出时，$\mathbb P_\sigma$ 恰好是环面簇 $\mathbb P_\Delta$；而 $\mathbb C[\sigma\cap\overline M]$ 就是\hyperref[subsec:toric-from-polytopes]{第 3.2.2 小节}的环 $S_\Delta$。超曲面 $Y\subset\mathbb P_\sigma$ 恰好是\hyperref[sec:batyrev-mirror]{第 4.1 节}研究的 Calabi--Yau 超曲面 $V\subset\mathbb P_\Delta$。
\item 当 $\sigma$ 由式~\eqref{eq:nef-gorenstein-cone} 给出时，存在双有理保典范态射
\[
\mathbb P\bigl(\mathcal O(E_1)\oplus\cdots\oplus\mathcal O(E_k)\bigr)
\longrightarrow\mathbb P_\sigma,
\]
使 $Y\subset\mathbb P_\sigma$ 是 $Y'\subset\mathbb P(\mathcal O(E_1)\oplus\cdots\oplus\mathcal O(E_k))$ 的像。
\end{itemize}
此外，在这两种情形中，镜像对称都化为 Gorenstein 锥的对偶性。

{\setlength{\emergencystretch}{1em}
除了这些已知例子，Gorenstein 锥构造也有助于解释通常镜像对称表述失效的一些情形。\emph{刚性 Calabi--Yau 三维簇}的情形最为明显。它们没有非平凡形变，所以 $H^{2,1}(V)=0$。通常意义下的镜像 $V^\circ$ 应有 $H^{1,1}(V^\circ)=0$，但这对 Kähler 流形是不可能的。然而，在某些情形中，人们仍找到了“镜像”。例如，设 $E$ 为唯一具有阶为 $3$ 的自同构的椭圆曲线，令 $V=E^3/G$，其中 $G$ 由该自同构的对角作用生成。可以证明 $V$ 是一个有奇点的刚性 Calabi--Yau 簇；\SourceCite{CDP} 算出的镜像是 $\mathbb P^8$ 中七维 Fermat 超曲面的一个有限商。\SourceCite{BB3} 表明，这与 Gorenstein 锥的对偶性一致。\par}\index{刚性Calabi--Yau三维簇@刚性 Calabi--Yau 三维簇|hyperpage}

\section{Voisin--Borcea 构造}\label{sec:voisin-borcea}

本节描述 Voisin 与 Borcea 各自独立给出的另一种镜像族构造。其灵感来自 Nikulin 关于带对合的 K3 曲面的工作。他找到了 $75$ 个这样的曲面族，其不变量显示出醒目的部分对称性 \SourceCite{Nikulin}。利用这些曲面，Voisin 与 Borcea 构造了猜想具有镜像对称性的 Calabi--Yau 三维簇族。我们将勾勒一些证明，完整细节请参见原始论文 \SourceCite{Voisin1,Borcea2}。\index{Voisin--Borcea构造@Voisin--Borcea 构造|hyperpage}

\SourcePage{89}
\subsection{由 K3 曲面与椭圆曲线构造三维簇}\label{subsec:voisin-threefolds}

首先利用某些带对合的 K3 曲面与椭圆曲线构造 Calabi--Yau 三维簇。更精确地说，考虑以下资料：
\begin{itemize}
\item 一个 K3 曲面 $S$ 与对合 $i:S\to S$，它在 $H^{2,0}(S)$ 上作用为 $-1$。$i$ 的不动点集是光滑曲线 $C_1,\ldots,C_N$ 的不交并。我们明确假设 $N>0$。
\item 一条椭圆曲线 $E$ 与对合 $j:E\to E$，其商为 $\mathbb P^1$。将 $j$ 的四个不动点记为 $p_1,p_2,p_3,p_4$。
\end{itemize}
$S$ 的对合 $i$ 与 $E$ 的对合 $j$ 诱导 $E\times S$ 的对合 $k$，定义为 $k(e,s)=(j(e),i(s))$，其中 $(e,s)\in E\times S$。$k$ 的不动点集由 $4N$ 条曲线 $C_{r,s}=\{p_r\}\times C_s$ 组成。沿这 $4N$ 条曲线爆破 $E\times S$，得到三维簇 $\widetilde{E\times S}$，其光滑例外除子为 $D=\bigcup D_{r,s}$；对合 $k$ 自然提升为 $\widetilde{E\times S}$ 的对合 $\widetilde k$。置
\begin{equation}\label{eq:voisin-threefold}
V=\widetilde{E\times S}/\widetilde k.
\end{equation}
可将 $D=\bigcup D_{r,s}$ 视为 $V$ 上具有 $4N$ 个不可约分支的除子。于是可以证明以下结论。
\begin{lemma}\label{lem:voisin-cy}
$V$ 是光滑 Calabi--Yau 三维簇。
\end{lemma}

令 $g_s$ 表示曲线 $C_s\subset S$ 的亏格，置 $N'=\sum_{s=1}^N g_s$。$V$ 的 Hodge 分量计算如下。
\begin{theorem}\label{thm:voisin-hodge}
$h^{1,1}(V)=11+5N-N'$，且 $h^{2,1}(V)=11+5N'-N$。
\end{theorem}
\begin{note}
由这些公式可见，若同样构造一个 $V^\circ$，但交换 $N$ 与 $N'$ 的角色，则它便是镜像流形的一个候选者。这正是 Voisin 与 Borcea 所做的。
\end{note}
\begin{proof}
先给出 $H^{1,1}(V)$ 的显式分解。首先，除子 $D=\bigcup D_{r,s}$ 的各分支在 $H^{1,1}(V)$ 中给出 $4N$ 个类。其次，光滑曲面 $T=S/i$ 给出一些类。将每条曲线 $C_s\subset S$ 与其在 $T$ 中的像等同，通过计算 Euler 示性数来求 $h^{1,1}(T)$：
\begin{equation}\label{eq:voisin-euler-quotient}
\begin{aligned}
e(T)&=e\left(T-\bigcup_s C_s\right)+\sum_s e(C_s)\\
&=\tfrac12e\left(S-\bigcup_s C_s\right)+\sum_s e(C_s)\\
&=\tfrac12\left(24-\sum_s(2-2g_s)\right)+\sum_s(2-2g_s)\\
&=12+N-N'.
\end{aligned}
\end{equation}
以 $i^*$ 表示 $i$ 诱导的 $S$ 的上同调上的作用。考察 $S$ 的上同调中 $i^*$ 不变的部分，立即可得 $T$ 的上同调满足 $h^0(T)=h^4(T)=1$，以及 $h^1(T)=h^3(T)=h^{2,0}(T)=h^{0,2}(T)=0$。结合式~\eqref{eq:voisin-euler-quotient} 得到
\begin{equation}\label{eq:voisin-quotient-h11}
h^{1,1}(T)=10+N-N'.
\end{equation}
$H^{1,1}(T)$ 中的上同调类可经自然投影 $\pi_1:V\to T$ 拉回到 $V$；将所得 $H^{1,1}(V)$ 中的元素与 $H^{1,1}(T)$ 等同。

\SourcePage{90}
最后，$H^{1,1}(E)$ 中的元素经自然投影 $\widetilde{E\times S}\to E$，拉回为 $H^{1,1}(\widetilde{E\times S})$ 中的 $\widetilde k$ 不变上同调类，因此也可与 $V$ 上的上同调类等同。

刚才描述的三类上同调类在 $\mathbb C$ 上线性无关，并张成 $H^{1,1}(V)$。因此得到直和分解
\begin{equation}\label{eq:voisin-h11-decomposition}
H^{1,1}(V)=\bigoplus_{r,s}\mathbb C\cdot[D_{r,s}]
\oplus H^{1,1}(T)\oplus H^{1,1}(E).
\end{equation}
此分解意味着 $h^{1,1}(V)=4N+(10+N-N')+1=11+5N-N'$，证明了所需的 $h^{1,1}(V)$ 公式。

对定理的第二个断言，将 $H^{2,1}(V)$ 与 $H^{2,1}(\widetilde{E\times S})$ 的 $\widetilde k$ 不变部分等同，再次列出三类上同调类。首先，由构造，$D_{r,s}$ 是 $C_{r,s}$ 上的射影丛，因此给出 $\widetilde{E\times S}$ 中以 $C_{r,s}$ 为参数的一族曲线。有相应的 Abel--Jacobi 映射 $H^{1,0}(C_{r,s})\to H^{2,1}(\widetilde{E\times S})$，见 \SourceCite{CG}；将其像与 $H^{1,0}(C_{r,s})$ 等同。其次，令 $H^{1,1}(S)^-$ 表示 $i^*$ 在 $H^{1,1}(S)$ 上作用的 $-1$ 特征子空间。由于 $j$ 在 $H^1(E)$ 上诱导 $-1$，而 $i$ 在 $H^{2,0}(S)$ 上作用为 $-1$，由 Künneth 分解，还得到来自 $H^{1,1}(S)^-\otimes H^{1,0}(E)$ 与 $H^{2,0}(S)\otimes H^{0,1}(E)$ 的类。刚才描述的三类仍在 $\mathbb C$ 上线性无关，并张成 $H^{2,1}(V)$。注意 $H^{1,0}(E)$ 与 $H^{0,1}(E)$ 都是一维的，因此可以非典范地写为
\begin{equation}\label{eq:voisin-h21-decomposition}
H^{2,1}(V)\simeq\bigoplus_{r,s}H^{1,0}(C_{r,s})
\oplus H^{1,1}(S)^-\oplus H^{2,0}(S).
\end{equation}
所以 $h^{2,1}(V)=4N'+h^{1,1}(S)^-+1$。不过，若 $H^{1,1}(S)^+$ 表示 $i^*$ 在 $H^{1,1}(S)$ 上作用的 $+1$ 特征子空间，则式~\eqref{eq:voisin-quotient-h11} 意味着
\begin{equation}\label{eq:voisin-eigenspace-h11}
h^{1,1}(S)^+=10+N-N',\qquad
h^{1,1}(S)^-=10+N'-N.
\end{equation}
所需的 $h^{2,1}(V)$ 公式立即随之得到。
\end{proof}

后面还需要对合 $i$ 的如下不变量。交叉形式 $\langle\cdot,\cdot\rangle$ 在 $H^2(S,\mathbb Z)$ 上是幺模的，但在 $H^2(S,\mathbb Z)^+$ 或 $H^2(S,\mathbb Z)^-$ 上未必如此。不过，利用 Nikulin 关于 K3 格的对合的结果 \SourceCite{AN,Dolgachev2,Nikulin}，可以证明存在整数 $a$，使
\begin{equation}\label{eq:voisin-discriminant}
\begin{aligned}
2^a&=\langle\cdot,\cdot\rangle\text{ 在 }H^2(S,\mathbb Z)^+
\text{ 上的判别式绝对值}\\
&=\langle\cdot,\cdot\rangle\text{ 在 }H^2(S,\mathbb Z)^-
\text{ 上的判别式绝对值}.
\end{aligned}
\end{equation}
$a$ 有如下显式公式。
\begin{lemma}\label{lem:voisin-discriminant}
$a=12-N-N'$。
\end{lemma}
\begin{proof}
关键步骤是构造正合列
\[
0\longrightarrow H^2(T,\mathbb Z)
\longrightarrow H^2(S,\mathbb Z)^+
\longrightarrow(\mathbb Z_2)^{N-1}\longrightarrow0,
\]
其中嵌入由拉回诱导。这里省略构造。$H^2(S,\mathbb Z)^+$ 上的形式 $\langle\cdot,\cdot\rangle$，限制到 $H^2(T,\mathbb Z)$ 后为其幺模形式的两倍，因此该正合列意味着判别式绝对值为 $2^{b_2(T)-2(N-1)}$。现在由式~\eqref{eq:voisin-quotient-h11} 即得引理。
\end{proof}

\SourcePage{91}
\subsection{Voisin--Borcea 镜像}\label{subsec:voisin-mirror}

现在描述 Calabi--Yau 三维簇 $V=\widetilde{E\times S}/\widetilde k$ 的镜像族 $V^\circ$。如\hyperref[subsec:voisin-threefolds]{第 4.4.1 小节}已指出，定理~\ref{thm:voisin-hodge} 表明，若存在同类型的镜像族，$N$ 与 $N'$ 的角色就必定互换。构造族 $V^\circ$ 需要以下技术性引理。
\begin{lemma}\label{lem:voisin-hyperbolic-plane}
若 $N'>0$ 且 $(N,N')\ne(5,1)$，则 $H^2(S,\mathbb Z)^-$ 含有一个双曲平面。
\end{lemma}
\SourceCite{Voisin1} 中的陈述实际上稍强一些。证明使用前面提到的 Nikulin 关于 K3 格的对合的结果。

假设 $N'>0$ 且 $(N,N')\ne(5,1)$，选择双曲平面 $P\subset H^2(S,\mathbb Z)^-$，以 $r_P$ 表示关于 $P$ 的反射。也就是说，$P$ 是 $r_P$ 的 $+1$ 特征子空间，$P^\perp$ 是其 $-1$ 特征子空间。置 $i_1^*=r_P\cdot i^*$。仍利用 Nikulin 关于 K3 格的对合的结果，可以证明 $i_1^*$ 是一个对合，它来自某个 K3 曲面 $S_1$ 上的对合 $i_1$。此外，简单计算表明
\begin{equation}\label{eq:voisin-mirror-lattice}
H^2(S_1,\mathbb Z)^-=P\oplus H^2(S,\mathbb Z)^+.
\end{equation}

现在由 $(S_1,i_1)$ 与带对合 $j_1$ 的椭圆曲线 $E_1$，按构造 $V$ 的相同方法构造 Calabi--Yau 三维簇 $V^\circ$。如 \SourceCite{Voisin1} 所述，$(E_1,j_1)$ 有一个自然选择。设 $V^\circ$ 的不变量为 $(N_1,N_1')$，则有以下结果。
\begin{proposition}\label{prop:voisin-mirror-invariants}
$N_1=N'$，$N_1'=N$。
\end{proposition}
\begin{proof}
设 $H^2(S_1,\mathbb Z)^-$ 上交叉形式的判别式绝对值为 $2^{a_1}$。因为 $P\subset H^2(S,\mathbb Z)$ 幺模，式~\eqref{eq:voisin-discriminant}、\eqref{eq:voisin-mirror-lattice} 意味着 $a_1=a$，所以由引理~\ref{lem:voisin-discriminant}，
\[
12-N_1-N_1'=12-N-N'.
\]
式~\eqref{eq:voisin-mirror-lattice} 还意味着 $h^{1,1}(S_1)^-=h^{1,1}(S)^+$。结合式~\eqref{eq:voisin-eigenspace-h11} 得到
\[
10+N_1'-N_1=10+N-N'.
\]
立即得到所需等式 $N_1=N'$ 与 $N_1'=N$。
\end{proof}

猜想所得 Calabi--Yau 三维簇族 $V^\circ$ 是 $V$ 的镜像族。作为支持 $V$ 与 $V^\circ$ 的镜像关系的证据，由定理~\ref{thm:voisin-hodge} 与命题~\ref{prop:voisin-mirror-invariants} 得到 $h^{1,1}(V)=h^{2,1}(V^\circ)$ 及 $h^{2,1}(V)=h^{1,1}(V^\circ)$。此外，Voisin 与 Borcea 还给出了其他支持镜像关系的证据。

第一项证据来自 Voisin，涉及 $V$ 与 $V^\circ$ 的复模和 Kähler 模之间的关系。$V$ 的复模可描述如下。先考虑来自 K3 曲面 $S$ 的模，考察区域
\[
D=\{\omega\in\mathbb P(H^2(S,\mathbb C)^-):
\omega\cdot\omega=0,\ \omega\cdot\overline\omega>0\}.
\]
它几乎是具有满足 $i'^*=i^*$ 的对合 $i'$ 的 K3 曲面的周期域。更精确地说，一般点 $\omega\in D$ 对应于带对合 $i'$ 的 K3 曲面 $S'$，满足 $i'^*=i^*$ 及 $\mathbb C\cdot\omega=H^{2,0}(S')$；有序对 $(S,i)$ 的模空间是某个明确的第二纲子集 $U\subset D$ 的离散商。固定椭圆曲线 $(E,j)$，便得到 $(S,i)$ 的模空间到 $V$ 的复模空间中的嵌入。因此，可将 $D$ 基本上视为比 $V$ 的复模空间某个叶状结构的一片叶稍大的空间。\index{对合|hyperpage}\index{双曲平面|hyperpage}

\SourcePage{92}
转向 $V$ 的 Kähler 模，定义区域
\[
D'=\{\alpha\in H^2(S,\mathbb C)^+:(\operatorname{Re}\alpha)^2>0\}.
\]
注意，$D'$ 可与 $H^2(T,\mathbb C)$ 的一个子集等同；粗略地说，它比 $T$ 的 Kähler 模空间稍大。此外，经 $\pi_1:V\to T$ 拉回后，$T$ 的 Kähler 模中的元素可加到 $V$ 的复化 Kähler 类上，得到新的复化 Kähler 类。因此，$T$ 的 Kähler 模空间可与 $V$ 的 Kähler 模空间某个叶状结构的各叶等同，而区域 $D'$ 比 $V$ 的 Kähler 模中的这片叶稍大。这与刚才对复模所作的处理很相似。

K3 曲面 $S$ 给出区域 $D$ 与 $D'$，它们分别与 $V$ 的复模和 Kähler 模相关。镜像一侧有上面构造的曲面 $S_1$，以 $D_1$、$D_1'$ 表示 $S_1$ 的相应区域。若 $V$ 与 $V^\circ$ 的确是镜像，则 $D$ 与 $D_1'$ 之间以及 $D'$ 与 $D_1$ 之间存在关系，将与镜像对称相一致。Voisin 证明了如下结果。
\begin{theorem}\label{thm:voisin-domains}
$D\simeq D_1'$，且 $D'\simeq D_1$。
\end{theorem}

Voisin 给出的另一类证据是 Yukawa 耦合的渐近行为。在 $H^{1,1}(V)$ 上有杯积给出的耦合；如\hyperref[chap:introduction]{第 1 章}定义，$H^1(V,T_V)$ 上的 Yukawa 耦合为
\begin{equation}\label{eq:voisin-yukawa}
\langle\theta_1,\theta_2,\theta_3\rangle
=\int_V\Omega\wedge
(\nabla_{\theta_1}\nabla_{\theta_2}\nabla_{\theta_3}\Omega),
\end{equation}
其中 $\Omega$ 是 $V$ 上的全纯三形式，见式~\eqref{eq:b-yukawa}。同构 $H^1(V,T_V)\simeq H^{2,1}(V)$ 在 $H^{2,1}(V)$ 上给出一个耦合，也称为 Yukawa 耦合。

利用式~\eqref{eq:voisin-h11-decomposition}、\eqref{eq:voisin-h21-decomposition} 的分解，Voisin 仔细描述了 $H^{1,1}(V)$ 与 $H^{2,1}(V^\circ)$ 上的耦合，并证明 $H^{1,1}(V)$ 上的杯积可渐近地与 $H^{2,1}(V^\circ)$ 上的 Yukawa 耦合等同。

为解释最后这个断言，回顾\hyperref[chap:introduction]{第 1 章}：$H^{2,1}(V^\circ)$ 上的 B 模型关联函数恰好就是上述 Yukawa 耦合，而 $H^{1,1}(V)$ 上的 A 模型关联函数由以下公式给出：
\[
\begin{aligned}
\langle\omega_1,\omega_2,\omega_3\rangle
={}&\int_V\omega_1\wedge\omega_2\wedge\omega_3\\
&+\sum_{\substack{\beta\ne0\\\beta\text{ 有效}}}n_\beta
\left(\int_\beta\omega_1\right)
\left(\int_\beta\omega_2\right)
\left(\int_\beta\omega_3\right)
\frac{e^{2\pi i\int_\beta\omega}}{1-e^{2\pi i\int_\beta\omega}},
\end{aligned}
\]
见式~\eqref{eq:a-yukawa}。现在设 $V$ 的复化 Kähler 类趋于无穷，即对所有有效曲线 $\beta$，都有 $\int_\beta\operatorname{Im}(\omega)\to\infty$；物理文献称之为“大半径极限”。于是
\[
\langle\omega_1,\omega_2,\omega_3\rangle
\longrightarrow\int_V\omega_1\wedge\omega_2\wedge\omega_3,
\]
右边就是杯积。在这种情形，称 A 模型关联函数渐近地变为杯积。

由于镜像流形 $V^\circ$ 的 B 模型关联函数与原流形 $V$ 的 A 模型关联函数一致，预期 B 模型关联函数也应具有类似行为。这正是 Voisin 证明的：在适当假设下，$V^\circ$ 上的 Yukawa 耦合有一个极限，与 $V$ 上的杯积一致，见 \SourceCite[第 3 节]{Voisin1}。

\SourcePage{93}
\subsection{与 Batyrev 构造的关系}\label{subsec:nonreflexive-weights}

Borcea 为 $V$ 与其 Voisin--Borcea 镜像 $V^\circ$ 的镜像对称提供了另一类证据：他证明，此构造的某些情形是 Batyrev 镜像对偶性的特殊情形。例如，考虑加权射影空间 $\mathbb P(3,3,2,2,2)$ 中的反典范超曲面
\begin{equation}\label{eq:voisin-weighted-original}
u^4+v^4=y^6+z^6+w^6,
\end{equation}
其中变量 $u,v,y,z,w$ 的权重依次为 $3,3,2,2,2$。为将它与式~\eqref{eq:voisin-threefold} 的三维簇联系起来，注意 $x^2=u^4+v^4$ 在 $\mathbb P(2,1,1)$ 中定义椭圆曲线 $E$，而 $x^2=y^6+z^6+w^6$ 在 $\mathbb P(3,1,1,1)$ 中定义（奇异的）K3 曲面 $S$。两种情形中，$x\mapsto-x$ 都是一个对合。Borcea 证明，Calabi--Yau 三维簇 $V=\widetilde{E\times S}/\widetilde k$ 是上述 $\mathbb P(3,3,2,2,2)$ 中超曲面的消解。还应注意，在 $S$ 中，对合的不动点集由 $y^6+z^6+w^6=0$ 给出，这是亏格为 $10$ 的光滑曲线。因此，此时 $(N,N')=(1,10)$。

利用 Voisin--Borcea 构造，族~\eqref{eq:voisin-weighted-original} 有镜像族 $V_{\mathrm{VB}}^\circ$，其不变量为 $(N,N')=(10,1)$。在 \SourceCite{Borcea2} 中，Borcea 证明，这个族包含加权射影空间 $\mathbb P(25,25,20,16,14)$ 中的反典范超曲面
\begin{equation}\label{eq:voisin-weighted-mirror}
u^4+v^4=y^5+yz^5+zw^6.
\end{equation}
此时，$E$ 仍是 $\mathbb P(2,1,1)$ 中的 $x^2=u^4+v^4$，但 $S_1$ 是 $\mathbb P(25,10,8,7)$ 中由 $x^2=y^5+yz^5+zw^6$ 定义的奇异 K3 曲面。如前，式~\eqref{eq:voisin-weighted-mirror} 的消解为 $V_{\mathrm{VB}}^\circ=\widetilde{E\times S_1}/\widetilde k$。对合的不动点集在 $S_1$ 中为 $y^5+yz^5+zw^6=0$。消解 $S_1$ 后，不动点集成为一条椭圆曲线加上九条 $\mathbb P^1$，这解释了为何 $(N,N')=(10,1)$。

基本断言是：$V_{\mathrm{VB}}^\circ$ 与式~\eqref{eq:voisin-weighted-original} 的 Batyrev 镜像一致，后者记为 $V_{\mathrm{Bat}}^\circ$。Borcea 在 \SourceCite{Borcea2} 中用\emph{\IndexLocation{indexloc:0075:93}双棱锥构造}证明了这一点；我们将利用 \SourceCite{CdK} 的\emph{广义 Berglund--Hübsch \IndexLocation{indexloc:0371:93}转置法则}给出更初等的证明。式~\eqref{eq:voisin-weighted-original} 与式~\eqref{eq:voisin-weighted-mirror} 之间的对偶性并不显然，因为 $\mathbb P(3,3,2,2,2)$ 是 Fano 环面簇，而 $\mathbb P(25,25,20,16,14)$ 不是，见引理~\ref{lem:weighted-projective-fano}。那么如何应用 Batyrev 对偶性？

答案涉及物理文献中的一种方法：把 Batyrev 镜像表示成非 Fano 加权射影空间中的超曲面。这里不描述一般情形，而说明它对 $\mathbb P(q)$、$q=(3,3,2,2,2)$ 如何起作用；更多细节见 \SourceCite[第 3.1 节]{CdK}。从环面几何角度看 $\mathbb P(q)$，有
\[
M=\{a\in\mathbb Z^5:q\cdot a=0\},\qquad
N=\mathbb Z^5/\mathbb Z q.
\]
$\mathbb P(q)$ 的扇的一维锥生成元为 $w_i=[e_i]\in N$，其中 $[e_i]$ 表示 $\mathbb Z^5$ 中第 $i$ 个标准基向量 $e_i$ 的陪集。此外，反典范除子的多胞形为
\[
\begin{aligned}
\Delta&=\{a\in M_{\mathbb R}:\langle a,w_i\rangle\geq-1\}\\
&=\{a\in M_{\mathbb R}:a\cdot e_i\geq-1\}.
\end{aligned}
\]
由于 $\mathbb P(q)$ 是 Fano 簇，$\Delta$ 自反。

\SourcePage{94}
现在考虑 $\Delta\cap M$ 中的以下各点：
\begin{equation}\label{eq:voisin-exponent-rows}
\begin{aligned}
a_0&=(3,-1,-1,-1,-1),\\
a_1&=(-1,3,-1,-1,-1),\\
a_2&=(-1,-1,4,0,-1),\\
a_3&=(-1,-1,-1,4,0),\\
a_4&=(-1,-1,-1,-1,5).
\end{aligned}
\end{equation}
这些点张成以 $0$ 为内点的单形。因此，某个系数全为正的 $a_i$ 的线性组合等于零，可以验证
\begin{equation}\label{eq:voisin-dual-weights}
25a_0+25a_1+20a_2+16a_3+14a_4=0.
\end{equation}
而且，其他所有关系都是这个关系的倍数。因此，$\widehat q=(25,25,20,16,14)$ 是各 $a_i$ 之间唯一的系数为正整数且互素的关系。它给出加权射影空间 $\mathbb P(\widehat q)=\mathbb P(25,25,20,16,14)$，具有
\[
\widehat M=\{a\in\mathbb Z^5:\widehat q\cdot a=0\},\qquad
\widehat N=\mathbb Z^5/\mathbb Z\widehat q.
\]
$\mathbb P(\widehat q)$ 的扇的一维锥生成元为 $\widehat w_i=[e_i]\in\widehat N$。

由于 $\Delta$ 自反，$\mathbb P(q)=\mathbb P_\Delta$ 是 Fano 簇；极多胞形 $\Delta^\circ$ 给出 Fano 环面簇 $\mathbb P_{\Delta^\circ}$，记为 $\mathbb P(q)^\circ$。$\mathbb P(q)^\circ$ 的扇位于 $M_{\mathbb R}$ 中，各锥由 $\Delta$ 的各面确定。为将 $\mathbb P(q)^\circ$ 与非 Fano 加权射影空间 $\mathbb P(\widehat q)$ 联系起来，使用映射 $\phi:\widehat N\to M$，将 $\widehat w_i\in\widehat N$ 送到 $a_i\in M$。由式~\eqref{eq:voisin-dual-weights}，容易看出 $\phi$ 良定义，并诱导同构 $\widehat N_{\mathbb R}\simeq M_{\mathbb R}$。因为 $\phi$ 把一维锥生成元 $\widehat w_i$ 送到 $a_i\in\Delta\cap M$，后者也可作为一维锥生成元，故能找到与 $\mathbb P(\widehat q)$ 的扇相容的极大射影细分 $\Sigma^\circ$。这意味着，将 $\Sigma^\circ$ 经 $\phi$ 拉回到 $\widehat N_{\mathbb R}$ 后所得的扇，细化 $\mathbb P(\widehat q)$ 的扇。由此得到环面爆破 $\widehat X\to\mathbb P(\widehat q)$，而 $\phi$ 诱导态射
\[
\phi_*:\widehat X\longrightarrow X_{\Sigma^\circ}.
\]
根据标准理论，$\phi_*$ 是有限态射，其次数为 $\phi(\widehat N)\subset M$ 的指数。令 $d=q\cdot(1,1,1,1,1)=12$、$\widehat d=\widehat q\cdot(1,1,1,1,1)=100$，则指数由下式计算：
\[
[M:\phi(\widehat N)]
=\frac{\det(A+\mathbf1)}{d\widehat d}
=\frac{2400}{12\cdot100}=2,
\]
其中 $A$ 是以式~\eqref{eq:voisin-exponent-rows} 中各向量为行的 $5\times5$ 矩阵，$\mathbf1$ 是所有元素均为 $1$ 的 $5\times5$ 矩阵。\SourceCite[第 3.1 节]{CdK} 陈述了这个公式但未给证明，不过论证并不困难。因此 $\phi_*:\widehat X\to X_{\Sigma^\circ}$ 是一个 $2:1$ 覆盖。

下一步是将 $V_{\mathrm{VB}}^\circ\subset\widehat X$ 与反典范超曲面 $V_{\mathrm{Bat}}^\circ\subset X_{\Sigma^\circ}$ 的拉回等同。自反单形 $\Delta^\circ\subset N_{\mathbb R}$ 的顶点为各 $w_i$，它们给出 Laurent 多项式 $f=\sum_{i=0}^4 c_it^{w_i}\in L(\Delta^\circ)$。它虽不是 $L(\Delta^\circ)$ 中最一般的 Laurent 多项式，却有正确的 Newton 多胞形。因此可设 $V_{\mathrm{Bat}}^\circ$ 由 $f=0$ 定义。为看出这在 $\widehat X$ 中给出什么，需要研究与 $\phi$ 对偶的映射 $\phi^t:N\to\widehat M$。将 $\mathbb Z^5$ 的元素视为列向量，并使用 $\widehat N$ 与 $M$ 的自然基及坐标，则复合映射
\[
\mathbb Z^5\longrightarrow\widehat N\xrightarrow{\phi}M
\longrightarrow\mathbb Z^5
\]
\SourcePage{95}
就是乘以 $A^t$，其中 $A$ 仍是以式~\eqref{eq:voisin-exponent-rows} 中各向量为行的矩阵。取对偶后，复合映射
\[
\mathbb Z^5\longrightarrow N\xrightarrow{\phi^t}\widehat M
\longrightarrow\mathbb Z^5
\]
就是乘以 $A$。这意味着 $\phi^t(w_i)=\widehat a_i$，其中 $\widehat a_i$ 是 $A$ 的第 $i$ 列。因此，$X_{\Sigma^\circ}$ 上的 $f=0$ 拉回为 $\widehat f=0$，其中 $\widehat f=\sum_{i=0}^4c_it^{\widehat a_i}$。最后一步是关于变量 $u,v,y,z,w$ 作齐次化。留给读者验证，这只需将 $\widehat a_i$ 的每个分量加 $1$。因此，得到的指数向量为矩阵
\[
A+\mathbf1=
\begin{pmatrix}
4&0&0&0&0\\
0&4&0&0&0\\
0&0&5&1&0\\
0&0&0&5&1\\
0&0&0&0&6
\end{pmatrix}
\]
的各列。用 $u,v,y,z,w$ 可将 $\widehat f=0$ 写为
\[
\widehat f=c_0u^4+c_1v^4+c_2y^5+c_3yz^5+c_4zw^6=0.
\]
式~\eqref{eq:voisin-weighted-mirror} 是这种形式的方程，因此可设 $V_{\mathrm{VB}}^\circ$ 由 $\widehat f=0$ 定义。

上述论证证明了 $\phi_*$ 诱导 $2:1$ 映射 $\phi_*:V_{\mathrm{VB}}^\circ\to V_{\mathrm{Bat}}^\circ$。这看起来不完全符合期望，因为我们想让二者相同。幸而有一个简单解释。在处理 $V_{\mathrm{VB}}^\circ=\widetilde{E\times S_1}/\widetilde k$ 时，将 $E$ 模去一个二阶点生成的子群，得到次数为二的同源映射 $E\to E'$，它与对合交换。因此有 $2:1$ 覆盖
\[
V_{\mathrm{VB}}^\circ\longrightarrow V_{\mathrm{VB}}'^\circ
=\widetilde{E'\times S_1}/\widetilde k.
\]
可以证明，$\phi_*:V_{\mathrm{VB}}^\circ\to V_{\mathrm{Bat}}^\circ$ 恰好是这种映射，所以可将 $V_{\mathrm{Bat}}^\circ$ 与 $V_{\mathrm{VB}}'^\circ$ 等同。证明使用 Borcea 的\IndexLocation{indexloc:0075:95}双棱锥构造，见 \SourceCite{Borcea2} 中式 (29) 后的说明。现在有一个关键观察：因为 $E$ 与 $E'$ 互为形变，$V_{\mathrm{VB}}^\circ$ 与 $V_{\mathrm{VB}}'^\circ$ 也互为形变。因此，就\emph{镜像族}而言，式~\eqref{eq:voisin-weighted-original} 的 Batyrev 镜像就是式~\eqref{eq:voisin-weighted-mirror}。这表明，至少在 $(N,N')=(1,10)$ 的情形，Voisin--Borcea 构造与 Batyrev 镜像对偶性一致。

论文 \SourceCite{Borcea2} 还列出了一些可用 Batyrev 镜像对偶性解释 Voisin--Borcea 构造的 $(N,N')$。不过，并非每个 Calabi--Yau 簇都是环面超曲面，所以可能并非所有 $(N,N')$ 都能这样处理。尽管如此，两种差别很大的\IndexLocation{indexloc:0799:95}\IndexLocation{indexloc:0804:95}镜像构造，即 Batyrev 构造与 Voisin--Borcea 构造，在式~\eqref{eq:voisin-weighted-original} 这样的情形中一致，是镜像对称确实存在的有力迹象。另一方面，两种构造的差异也表明，我们离一个真正普遍适用的寻找镜像的方法还有很远。

最后谈谈\hyperref[subsec:batyrev-hodge]{第 4.1.3 小节}提到的 $7555$ 个加权射影空间的列表。与 $\mathbb P(25,25,20,16,14)$ 一样，表中大部分加权射影空间\emph{不是} Fano 的。列表只显示出部分对称性，人们因而担心那些“缺失的镜像”。这个困难的解决使用了 Batyrev 镜像对偶性，以及上述用非 Fano 加权射影空间解释 Batyrev 对偶的方法的一个变体。从这个角度看，“缺失的镜像”对应于那些不能用加权射影空间表示的 Batyrev 对偶。细节与例子参见 \SourceCite{CdK}。

